% eksperymentalne tłumaczenie na włoski

%

\pol{\section{Nierówności trójkątne}}
\ita{\section{Disuguaglianze del triangolo}}

\begin{proposition}
% [stosunek sumy środkowych do obwodu] % lang-pl
% [stosunek sumy środkowych do obwodu] % lang-it
    Dany jest trójkąt o bokach $a, b, c$ oraz środkowych $m_a, m_b, m_c$.
    Wtedy
    \begin{equation}
        \frac 3 4 < \frac{m_a + m_b + m_c}{a + b + c} < 1.
    \end{equation}
\end{proposition}

Joanna Jaszuńska poda w $\Delta_{16}^{6}$ siedem nierówności dla długości boków trójkąta, do których przydaje się podstawienie Raviego $a = y + z$, $b = z + x$, $c = x + y$ (liczby $x, y, z > 0$ są długościami odcinków stycznych do okręgu wpisanego w trójkąt); wśród nich będzie nierówność $abc \ge (-a + b + c)(a - b + c)(a + b - c)$, której autorem jest Alessandro Padoa, oraz nierówność Alfreda Nesbitta. % lang-pl
Joanna Jaszuńska presenterà in $\Delta_{16}^{6}$ sette disuguaglianze per le lunghezze dei lati di un triangolo, per le quali è utile la sostituzione di Ravi $a = y + z$, $b = z + x$, $c = x + y$ (i numeri $x, y, z > 0$ sono le lunghezze dei segmenti di tangenza alla circonferenza inscritta nel triangolo); tra di esse ci saranno la disuguaglianza $abc \ge (-a + b + c)(a - b + c)(a + b - c)$, il cui autore è Alessandro Padoa, e la disuguaglianza di Alfred Nesbitt. % lang-it

\todofoot{Coxeter, s. 12}
\todofoot{https://www.deltami.edu.pl/2016/06/boki-trojkata/}

%

\subsection{Nierówność izoperymetryczna} % lang-pl
\subsection{Disuguaglianza isoperimetrica} % lang-it

\begin{proposition}
[nierówność izoperymetryczna] % lang-pl
[disuguaglianza isoperimetrica] % lang-it
    Dany jest trójkąt o obwodzie $2p$ oraz polu $S$. % lang-pl
    Dato un triangolo con perimetro $2p$ e area $S$. % lang-it
    Wtedy % lang-pl
    Allora % lang-it
    \begin{equation}
        S \le \frac{p^2}{3 \sqrt 3},
    \end{equation}
    zatem wśród trójkątów o ustalonym obwodzie największe pole ma trójkąt równoboczny. % lang-pl
    pertanto, tra i triangoli con perimetro fissato, quello equilatero ha l'area massima. % lang-it
\index{nierówność!izoperymetryczna}% % lang-pl
\index{disuguaglianza!isoperimetrica}% % lang-it
    % TODO: https://www.deltami.edu.pl/2011/04/symetryzacja-steinera/
\end{proposition}

Guzicki \cite[s. 169, 170]{guzicki_2021} wyprowadza nierówność izoperymetryczną z nierówności między średnimi oraz ze wzoru Herona, a później \cite[s. 362]{guzicki_2021} raz jeszcze, trygonometrycznie. % lang-pl
Guzicki \cite[p. 169, 170]{guzicki_2021} ricava la disuguaglianza isoperimetrica dalle disuguaglianze tra medie e dalla formula di Erone, e successivamente \cite[p. 362]{guzicki_2021} di nuovo, trigonometricamente. % lang-it
\index{wzór!Herona}% % lang-pl
\index{formula!di Erone}% % lang-it
Geometryczny dowód nierówności między średnimi harmoniczną, geometryczną, arytmetyczną i kwadratową dwóch liczb dodatnich $x$, $y$ poda Jarosław Górnicki w $\Delta_{91}^{7}$: na odcinku $AC$ długości $x + y$, z punktem $B$ takim, że $|AB| = x$ i $|BC| = y$, rysuje półokrąg i porównuje odpowiednie odcinki. Zadanie M 601 z $\Delta_{91}^{9}$ orzeka, że pole wielokąta o średnicy 1 nie przekracza $\pi/4$. % lang-pl
Una dimostrazione geometrica della disuguaglianza tra le medie armonica, geometrica, aritmetica e quadratica di due numeri positivi $x$, $y$ sarà fornita da Jarosław Górnicki in $\Delta_{91}^{7}$: sul segmento $AC$ di lunghezza $x + y$, con il punto $B$ tale che $|AB| = x$ e $|BC| = y$, egli traccia una semicirconferenza e confronta i segmenti opportuni. Il problema M 601 di $\Delta_{91}^{9}$ afferma che l'area di un poligono di diametro 1 non supera $\pi/4$. % lang-it
Bartłomiej Bzdęga opisze w $\Delta_{23}^1$ nierówność $4 \pi S \le L^2$, gdzie $L$ oznacza obwód figury, $S$ jej pole, zaś równość ma miejsce dokładnie dla koła. % lang-pl
Bartłomiej Bzdęga descriverà in $\Delta_{23}^1$ la disuguaglianza $4 \pi S \le L^2$, dove $L$ indica il perimetro della figura, $S$ la sua area, e l'uguaglianza si verifica esattamente per il cerchio. % lang-it
(Mamy też nierówność Bonnesena, która orzeka, że $\pi (R-r)^2 \le 4p^2 - 4 \pi L$). % lang-pl
(Esiste anche la disuguaglianza di Bonnesen, che stabilisce che $\pi (R-r)^2 \le 4p^2 - 4 \pi L$). % lang-it

Odpowiednik tego w trzech wymiarach znany będzie jako hipoteza (Zdzisława Aleksandra) Melzaka: podana około roku 1965, będzie musiała długo czekać na swoje rozwiązanie. % lang-pl
L'equivalente in tre dimensioni sarà noto come ipotesi di (Zdzisław Aleksander) Melzak: proposta intorno al 1965, dovrà attendere a lungo la sua soluzione. % lang-it
\index{hipoteza Melzaka}% % lang-pl
\index{ipotesi di Melzak}% % lang-it
\index[persons]{Melzak, Zdzisław Aleksander}%
% https://math.stackexchange.com/questions/4044670/what-is-the-largest-volume-of-a-polyhedron-whose-skeleton-has-total-length-1-is
Podamy trzy przykłady brył o obwodzie (sumie krawędzi) równym $1$: czworościan foremny, sześcian i graniastosłup trójkątny (prawdopodobnie maksymalizujący objętość). % lang-pl
Mostreremo tre esempi di solidi con perimetro (somma degli spigoli) pari a $1$: tetraedro regolare, cubo e prisma triangolare (probabilmente massimizzante il volume). % lang-it
Ich objętości wynoszą kolejno: % lang-pl
I loro volumi sono rispettivamente: % lang-it
\begin{align}
    V_4 & = \frac{\sqrt{2}}{2592} \approx 0.00054\,56070, \\
    V_6 & = \frac{1}{1728} \approx 0.00057\,87037, \\
    V_9 & = \frac{\sqrt{3}}{2916} \approx 0.00059\,39817.
\end{align}

%
%

\subsection{Nierówność Eulera} % lang-pl
\subsection{Disuguaglianza di Eulero} % lang-it

\begin{proposition}
[nierówność Eulera] % lang-pl
[disuguaglianza di Eulero] % lang-it
\index{nierówność!Eulera}% % lang-pl
\index{disuguaglianza!di Eulero}% % lang-it
	Niech $R, r, d$ oznaczają kolejno promienie okręgu opisanego i wpisanego w~trójkąt oraz odległość między środkami tych okręgów. % lang-pl
	Wtedy % lang-pl
	Siano $R, r, d$ rispettivamente i raggi della circonferenza circoscritta e inscritta a un triangolo e la distanza tra i centri di tali circonferenze. % lang-it
	Allora % lang-it
	\begin{equation}
		\frac{1}{R-d} + \frac{1}{R+d} = \frac 1 r,
	\end{equation}
	lub równoważnie $d^2 = R \cdot (R - 2r)$. % lang-pl
	oppure, equivalentemente, $d^2 = R \cdot (R - 2r)$. % lang-it
\end{proposition}

Leonhard Euler opublikuje ten wynik w 1765 roku, chociaż wcześniej zrobi to William Chapple, w pracy \emph{,,An essay on the properties of triangles inscribed in and circumscribed about two given circles''} z 1746 roku, na dole strony 123. % lang-pl
Leonhard Euler pubblicò questo risultato nel 1765, sebbene William Chapple lo avesse già pubblicato in precedenza, nel lavoro \emph{``An essay on the properties of triangles inscribed in and circumscribed about two given circles''} del 1746, in fondo a pagina 123. % lang-it
% Leversha, Gerry; Smith, G. C. (November 2007), "Euler and triangle geometry", The Mathematical Gazette, 91 (522): 436–452, doi:10.1017/S0025557200182087, JSTOR 40378417, S2CID 125341434
% TODO: https://en.wikipedia.org/wiki/Euler%27s_theorem_in_geometry a stornger version... [6]
Wspomni go Zetel \cite[s. 75]{zetel_2020}, Guzicki \cite[s. 362, 368-370]{guzicki_2021}. % lang-pl
Lo menzionano Zetel \cite[p. 75]{zetel_2020} e Guzicki \cite[pp. 362, 368-370]{guzicki_2021}. % lang-it
W rubryce Klub 44, redagowanej przez Marcina Kuczmę, pojawi się w $\Delta_{84}^{1}$ rozwiązanie zadania o znaku wyrażenia $2R + r - p$, gdzie $p$ oznacza połowę obwodu trójkąta: okaże się, że zależy on od tego, czy trójkąt jest ostrokątny, prostokątny czy rozwartokątny. % lang-pl
Nella rubrica Klub 44, curata da Marcin Kuczma, comparirà in $\Delta_{84}^{1}$ la soluzione di un problema sul segno dell'espressione $2R + r - p$, dove $p$ indica il semiperimetro del triangolo: risulterà che esso dipende dal fatto che il triangolo sia acutangolo, rettangolo oppure ottusangolo. % lang-it

\begin{corollary}
	Wynika stąd, że $R \ge 2r$. % lang-pl
	Ne segue che $R \ge 2r$. % lang-it
\end{corollary}

Trójwymiarowym odpowiednikiem jest: % lang-pl
L'analogo tridimensionale è il seguente: % lang-it

\begin{proposition}
[nierówność Grace'a-Danielssona] % lang-pl
[disuguaglianza di Grace-Danielsson] % lang-it
\index{nierówność!Grace'a-Danielssona}% % lang-pl
\index{disuguaglianza!di Grace-Danielsson}% % lang-it
	Dla każdej pary sfer o promieniach $r, R$ takich, że $r < R$ istnieje czworościan zawarty w dużej sferze i zawierający małą sferę wtedy i tylko wtedy, gdy % lang-pl
	Per ogni coppia di sfere di raggi $r, R$ tali che $r < R$ esiste un tetraedro inscritto nella sfera maggiore e circoscritto alla sfera minore se e solo se % lang-it
	\begin{equation}
		d^2 \le (R+r)(R - 3r).
	\end{equation}
\end{proposition}

% TODO: Grace, J.H. (1918), Proc. London Math. (ed.), Tetrahedra in relation to spheres and quadrics, vol. Soc.17, pp. 259–271
% Danielsson, G. (1952), Johan Grundt Tanums Forlag (ed.), Proof of the inequality d2≤(R+r)(R−3r) for the distance between the centres of the circumscribed and inscribed spheres of a tetrahedron, pp. 101–105

%
%

Wzmocnienie nierówności Eulera: % lang-pl
Un rafforzamento della disuguaglianza di Eulero: % lang-it

\begin{proposition}%
[nierówność Mitrinovica] % lang-pl
[disuguaglianza di Mitrinović] % lang-it
\index{nierówność!Mitrinovica}% % lang-pl
\index{disuguaglianza!di Mitrinović}% % lang-it
    Dany jest trójkąt o bokach $a, b, c$, wpisany w okrąg o promieniu $R$, opisany na okręgu o promieniu $r$. % lang-pl
    Dato un triangolo di lati $a$, $b$, $c$, inscritto in una circonferenza di raggio $R$ e circoscritto a una circonferenza di raggio $r$. % lang-it
    Wtedy % lang-pl
    Sia % lang-it
    % https://www.cut-the-knot.org/m/Geometry/InequalitiesInTriangle.shtml
    \begin{equation}
        6 \sqrt{3} r \le a + b + c \le 3 \sqrt{3} R.
    \end{equation}
\end{proposition}

Guzicki \cite[s. 362, 363]{guzicki_2021} dowodzi tego z nierówności izoperymetrycznej oraz oszacowania % lang-pl
Guzicki \cite[p. 362, 363]{guzicki_2021} dimostra questo risultato a partire dalla disuguaglianza isoperimetrica e dalla stima % lang-it
\begin{equation}
    \sin \alpha + \sin \beta + \sin \gamma \le 3 \sqrt{3} / 2.
\end{equation}

Marcin Pilipczuk wywnioskuje tę nierówność w $\Delta_{08}^{6}$ z nierówności Jensena dla funkcji wypukłej $-\sin x$ na przedziale $[0, \pi]$ i zauważy, że spośród trójkątów wpisanych w okrąg o promieniu $1$ największy obwód ma trójkąt równoboczny (z twierdzenia sinusów obwód wynosi $2 \sin \alpha + 2 \sin \beta + 2 \sin \gamma$). % lang-pl
Marcin Pilipczuk dedurrà questa disuguaglianza in $\Delta_{08}^{6}$ dalla disuguaglianza di Jensen per la funzione convessa $-\sin x$ sull'intervallo $[0, \pi]$ e osserverà che, tra i triangoli inscritti in una circonferenza di raggio $1$, quello di perimetro massimo è il triangolo equilatero (per il teorema dei seni il perimetro è $2 \sin \alpha + 2 \sin \beta + 2 \sin \gamma$). % lang-it

% TODO: https://www.cut-the-knot.org/m/Geometry/InequalitiesInTriangle.shtml

%
% eksperymentalne tłumaczenie na włoski

%

\subsection{Nierówność Leibniza} % lang-pl
\subsection{Disuguaglianza di Leibniz} % lang-it

U Zetela \cite[s. 72]{zetel_2020} znajdziemy wniosek z twierdzenia Stewarta (\ref{twierdzenie_stewarta}): % lang-pl
In Zetel \cite[p. 72]{zetel_2020} troviamo una conseguenza del teorema di Stewart (\ref{twierdzenie_stewarta}): % lang-it

\begin{proposition}
    Odległość między środkiem ciężkości i środkiem okręgu opisanego na trójkącie wynosi % lang-pl
    La distanza tra il baricentro e il circocentro di un triangolo è % lang-it
    \begin{equation}
        \sqrt{R^2 - \frac 19 \left(a^2 + b^2 + c^2\right)},
    \end{equation}
\end{proposition}

Wynika z niego: % lang-pl
Ne consegue: % lang-it

\begin{proposition}
[nierówność Leibniza] % lang-pl
[disuguaglianza di Leibniz] % lang-it
    Niech... % lang-pl
    Sia... % lang-it
    \begin{equation}
        a^2 + b^2 + c^2 \le 9R^2,
    \end{equation}
    z równością dla trójkąta równobocznego. % lang-pl
    \index{nierówność!Leibniza} % lang-pl
    con uguaglianza nel caso del triangolo equilatero. % lang-it
    \index{disuguaglianza!di Leibniz} % lang-it
\end{proposition}

Guzicki \cite[s. 362, 363, 371-375]{guzicki_2021} dowodzi tego trygonometrycznie: % lang-pl
Guzicki \cite[p. 362, 363, 371-375]{guzicki_2021} dimostra questo risultato mediante la trigonometria: % lang-it
\begin{equation}
    \sin^2 \alpha + \sin^2 \beta + \sin^2 \gamma = 2 (1 + \cos \alpha \cos \beta \cos \gamma) \le \frac 94.
\end{equation}

Twierdzenie Leibniza w postaci: suma kwadratów odległości dowolnego punktu $P$ od wierzchołków trójkąta jest równa sumie kwadratów odległości wierzchołków od środka ciężkości powiększonej o potrojony kwadrat odległości $P$ od środka ciężkości, wyprowadzi z własności momentu bezwładności (twierdzenia Steinera) Radosław Szmytkowski w $\Delta_{87}^{1}$; w ćwiczeniach poprosi też o wyznaczenie odległości środka okręgu opisanego od środka ciężkości trójkąta. % lang-pl
Il teorema di Leibniz nella forma: la somma dei quadrati delle distanze di un punto qualsiasi $P$ dai vertici di un triangolo è uguale alla somma dei quadrati delle distanze dei vertici dal baricentro, aumentata del triplo del quadrato della distanza di $P$ dal baricentro, sarà dedotto dalle proprietà del momento d'inerzia (teorema di Steiner) da Radosław Szmytkowski in $\Delta_{87}^{1}$; tra gli esercizi chiederà anche di determinare la distanza del circocentro dal baricentro del triangolo. % lang-it

%
% eksperymentalne tłumaczenie na włoski

%

\subsection{Nierówność Weitzenböcka i jej uogólnienia} % lang-pl

\subsection{Disuguaglianza di Weitzenböck e sue generalizzazioni} % lang-it

% https://en.wikipedia.org/wiki/Weitzenböck's_inequality

Roland Weitzenböck \cite{weitzenbock_1919} udowodni w 1919 roku: % lang-pl

Nel 1919 Roland Weitzenböck \cite{weitzenbock_1919} dimostrerà: % lang-it

\index[persons]{Weitzenböck, Roland}%

\begin{proposition}[nierówność Weitzenböcka] % lang-pl
	W trójkącie o bokach długości $a, b, c$ oraz polu powierzchni $S$ zachodzi nierówność % lang-pl
\begin{proposition}[disuguaglianza di Weitzenböck] % lang-it
    In un triangolo con lati di lunghezza $a$, $b$, $c$ e area $S$ vale la disuguaglianza % lang-it
    \begin{equation}
        a^2 + b^2 + c^2 \ge 4 \sqrt{3} S,
    \end{equation}
    z równością wtedy i tylko wtedy, gdy trójkąt jest równoboczny: $a = b = c$. % lang-pl
    
	\index{nierówność!Weitzenböcka} % lang-pl
    con uguaglianza se e solo se il triangolo è equilatero: $a = b = c$. % lang-it
    \index{disuguaglianza!di Weitzenböck} % lang-it
\end{proposition}

% https://en.wikipedia.org/wiki/Weitzenböck%27s_inequality#Further_proofs

Dowód wykorzystuje punkt Fermata, wzór Herona, nierówność między średnią arytmetyczną i~geometryczną, własności funkcji cotangens albo nierówność % lang-pl

La dimostrazione utilizza il punto di Fermat, la formula di Erone, la disuguaglianza tra media aritmetica e geometrica, le proprietà della funzione cotangente oppure la disuguaglianza % lang-it

\begin{equation}
    (a^2 - b^2)^2 + (b^2 - c^2)^2 + (c^2 - a^2)^2 \ge 0,
\end{equation}

patrz na przykład do Guzickiego \cite[s. 363, 365-367]{guzicki_2021}. % lang-pl
si veda ad esempio Guzicki \cite[pp.~363, 365--367]{guzicki_2021}. % lang-it


\index{punkt Fermata} % lang-pl
\index{punto di Fermat} % lang-it


\index{wzór Herona} % lang-pl
\index{formula!di Erone} % lang-it

% https://www.ime.usp.br/~toscano/disc/2022/GreenbergGeometry.pdf ps. 34

Nierówność pojawi się też w niezmienionej postaci na III Międzynarodowej Olimpiadzie Matematycznej. % lang-pl
La disuguaglianza comparirà inoltre, nella stessa forma, alla III Olimpiade Internazionale della Matematica. % lang-it


Są dwa uogólnienia: nierówność Hadwigera-Finslera oraz nierówność (Neuberga-)Pedoego. % lang-pl
Esistono due generalizzazioni: la disuguaglianza di Hadwiger--Finsler e la disuguaglianza di Neuberg--Pedoe. % lang-it

\begin{proposition}[nierówność Hadwigera-Finslera] % lang-pl
    W trójkącie o bokach długości $a, b, c$ oraz polu powierzchni $S$ zachodzi nierówność % lang-pl
\begin{proposition}[disuguaglianza di Hadwiger--Finsler] % lang-it
    
    In un triangolo con lati di lunghezza $a$, $b$, $c$ e area $S$ vale la disuguaglianza % lang-it
    \begin{equation}
        a^2 + b^2 + c^2 \ge (a-b)^2 + (b-c)^2 + (a-c)^2 + 4 \sqrt{3} S,
    \end{equation}
\end{proposition}

\begin{proof}
    Z twierdzenia kosinusów mamy $a^2 = b^2 + c^2 - 2bc \cos \alpha$. % lang-pl
    
    Dal teorema del coseno otteniamo $a^2 = b^2 + c^2 - 2bc \cos \alpha$. % lang-it

    Ponieważ $S = \frac 1 2 b c \sin \alpha$, poprzednią równość możemy przepisać do % lang-pl
    
    Poiché $S = \frac 1 2 b c \sin \alpha$, possiamo riscrivere l'uguaglianza precedente come % lang-it
    \begin{align}
        a^2 & = (b-c)^2 + 4S \cdot \frac{1 - \cos \alpha}{\sin \alpha} \\
        & = (b-c)^2 + 4S \cdot \frac{2 \sin^2 (\alpha / 2)}{2 \sin (\alpha / 2) \cos (\alpha / 2)} \\
        & = (b-c)^2 + 4S \cdot \tan \frac \alpha 2.
    \end{align}

    Po powtórzeniu tego dla pozostałych boków i dodaniu do siebie dostaniemy % lang-pl
    
    Ripetendo lo stesso ragionamento per gli altri lati e sommando le uguaglianze otteniamo % lang-it
    \begin{equation}
        a^2 + b^2 + c^2 = (a-b)^2 + (b-c)^2 + (c-a)^2 + 4S \left(\tan \frac \alpha 2 + \tan \frac \beta 2 + \tan \frac \gamma 2\right).
    \end{equation}

    Ale połówki kątów trójkąta nie przekraczają $\pi/2$, zaś funkcja $\tan$ jest wypukła, więc nawias można ograniczyć z dołu: % lang-pl
    
    Poiché le metà degli angoli del triangolo non superano $\pi/2$ e la funzione $\tan$ è convessa, possiamo stimare dal basso la parentesi: % lang-it
    \begin{equation}
        \tan \frac \alpha 2 + \tan \frac \beta 2 + \tan \frac \gamma 2 \ge 3 \tan \frac{\alpha + \beta + \gamma}{6} = 3\tan \frac \pi 6 = \sqrt 3,
    \end{equation}

    co kończy dowód. % lang-pl
    
    e questo conclude la dimostrazione. % lang-it
\end{proof}

Paul Finsler oraz Hugo Hadwiger \cite{finsler_1937} udowodnią to twierdzenie w 1937 roku. % lang-pl

Paul Finsler e Hugo Hadwiger \cite{finsler_1937} dimostreranno questo teorema nel 1937. % lang-it

\index[persons]{Finsler, Paul}%
\index[persons]{Hadwiger, Hugo}%

\begin{proposition}[nierówność Neuberga-Pedoego] % lang-pl
    W dwóch trójkątach o bokach $a, b, c$ oraz polu powierzchni $S$ i bokach $a_*, b_*, c_*$ oraz polu powierzchni $S_*$ zachodzi nierówność % lang-pl
\begin{proposition}[disuguaglianza di Neuberg--Pedoe] % lang-it
    
    Per due triangoli con lati $a$, $b$, $c$ e area $S$, e lati $a_*$, $b_*$, $c_*$ e area $S_*$, vale la disuguaglianza % lang-it
    \begin{equation}
        a^2_*(b^2 + c^2 - a^2) + 
        b^2_*(a^2 + c^2 - b^2) + 
        c^2_*(a^2 + b^2 - c^2) \ge 16 SS_*
    \end{equation}
    z równością wtedy i tylko wtedy, gdy trójkąty są podobne: $a/a_* = b/b_* = c/c_*$. % lang-pl
    
    con uguaglianza se e solo se i triangoli sono simili: $a/a_* = b/b_* = c/c_*$. % lang-it
\end{proposition}

(Aby dostać nierówność Weitzenböcka, wystarczy wziąć $a = b = c = 1$). % lang-pl
(Per ottenere la disuguaglianza di Weitzenböck, basta porre $a = b = c = 1$). % lang-it


Daniel Pedoe \cite{pedoe_1941} odkryje nierówność w 1941, a później dowie się, że była znana luksemburskiemu geometrze Józefowi Neubergowi. % lang-pl
Daniel Pedoe \cite{pedoe_1941} scoprirà la disuguaglianza nel 1941 e in seguito verrà a sapere che era già nota al geometra lussemburghese Joseph Neuberg. % lang-it

\index[persons]{Pedoe, Daniel}%
\index[persons]{Neuberg, Joseph}%

Ten jednak nie pokaże, że z równości dwóch strony wynika podobieństwo trójkątów. % lang-pl
Quest'ultimo, tuttavia, non mostrerà che l'uguaglianza tra i due membri implica la similitudine dei triangoli. % lang-it


Jeden z dowodów wykorzystuje wzór Herona oraz nierówność Cauchy'ego-Schwarza. % lang-pl
Una delle dimostrazioni utilizza la formula di Erone e la disuguaglianza di Cauchy--Schwarz. % lang-it

%
% eksperymentalne tłumaczenie na włoski

%

\subsection{Nierówność Ono} % lang-pl
\subsection{Disuguaglianza di Ono} % lang-it

% https://en.wikipedia.org/wiki/Ono%27s_inequality

W 1914 roku Toda Ono wysnuje hipotezę (nieprawdziwą w ogólności, ale prawdziwą dla wszystkich trójkątów ostrokątnych): % lang-pl
Nel 1914 Toda Ono proporrà la seguente congettura (falsa in generale, ma vera per tutti i triangoli acutangoli): % lang-it
\index[persons]{Ono, Toda}%

\begin{proposition}
[nierówność Ono] % lang-pl
[disuguaglianza di Ono] % lang-it
    W trójkącie ostrokątnym o bokach długości $a, b, c$ oraz polu powierzchni $S$ zachodzi nierówność % lang-pl
    In un triangolo acutangolo con lati di lunghezza $a$, $b$, $c$ e area $S$ vale la disuguaglianza % lang-it
    \begin{equation}
        \frac{27}{4096} \cdot (b^2 + c^2 - a^2)^2 (c^2 + a^2 - b^2)^2 (a^2 + b^2 - c^2)^2 \le S^6.
    \end{equation}
\end{proposition}

Dowód znajdzie F. Balitrand w 1916 roku. % lang-pl
La dimostrazione sarà trovata da F. Balitrand nel 1916. % lang-it
\index[persons]{Balitrand, F?}%
% Ono, T. (1914). "Problem 4417". Interméd. Math. 21: 146.
% https://zbmath.org/?format=complete&q=an:46.0859.06
Nierówność nie zachodzi dla trójkąta o bokach długości $a = 2$, $b = 3$, $c = 4$, wtedy $S^2 = 135/16$. % lang-pl
La disuguaglianza non vale per il triangolo di lati $a = 2$, $b = 3$, $c = 4$; in tal caso $S^2 = 135/16$. % lang-it

%
%

\index{nierówność!Erdősa-Mordella|(} % lang-pl
\index{disuguaglianza!di Erdős-Mordell|(} % lang-it

\subsection{Nierówność Erdősa-Mordella i jej uogólnienie} % lang-pl
\subsection{Disuguaglianza di Erdős-Mordell e sua generalizzazione} % lang-it

\begin{theorem}
[nierówność Erdősa-Mordella] % lang-pl
[disuguaglianza di Erdős-Mordell] % lang-it
    Niech $P$ będzie punktem wewnątrz trójkąta $\triangle ABC$, zaś $A_p, B_p, C_p$ spodkami punktu $P$ na boki trójkąta jak na rysunku \ref{erdos_mordell_barrowa}. % lang-pl
    Sia $P$ un punto interno al triangolo $\triangle ABC$, e siano $A_p, B_p, C_p$ le proiezioni ortogonali del punto $P$ sui lati del triangolo, come nella figura \ref{erdos_mordell_barrowa}. % lang-it
    Wtedy % lang-pl
    Allora % lang-it
    \begin{equation}
        |PA| + |PB| + |PC| \ge 2 (|PA_p| + |PB_p| + |PC_p|).
    \end{equation}
\end{theorem}

\begin{figure}[H] \centering
\begin{minipage}[b]{.45\linewidth}
\begin{center}
    \begin{comment}
    \begin{tikzpicture}[scale=.4]
    \tkzDefPoint(0, 0){A}
    \tkzDefPoint(10, 2){B}
    \tkzDefPoint(6, 7){C}
    \tkzDefPoint(5, 3){P}
    \tkzLabelPoint[below left](A){$A$}
    \tkzLabelPoint[below right](B){$B$}
    \tkzLabelPoint[above](C){$C$}
    \tkzLabelPoint[below left](P){$P$}
    \tkzDefPointsBy[projection=onto A--B](P){Pc}
    \tkzDefPointsBy[projection=onto B--C](P){Pa}
    \tkzDefPointsBy[projection=onto C--A](P){Pb}
    \tkzLabelPoint[above right](Pa){$A_p$}
    \tkzLabelPoint[above left](Pb){$B_p$}
    \tkzLabelPoint[below](Pc){$C_p$}

    \tkzDrawSegments[line width=0.2mm,dashed](P,Pa P,Pb P,Pc)
    \tkzDrawPolygon[line width=0.3mm](A,B,C)
    \tkzMarkRightAngles[size=0.5](P,Pa,C P,Pb,A P,Pc,B)
    \tkzDrawPoints[size=3,color=black,fill=black!50](A,B,C,P,Pc,Pb,Pa)
\end{tikzpicture}
\end{comment}
    \end{center}
    \subcaption{%
nierówność Erdősa-Mordella% % lang-pl
disuguaglianza di Erdős-Mordell% % lang-it
}
    \label{erdos_mordell_barrowa}
\end{minipage}
%
\begin{minipage}[b]{.45\linewidth}
\begin{center}\begin{comment}
    \begin{tikzpicture}[scale=.4]
    \tkzDefPoint(0, 0){A}
    \tkzDefPoint(10, 2){B}
    \tkzDefPoint(6, 7){C}
    \tkzDefPoint(5, 3){P}

    \tkzDefLine[bisector](A,P,B) \tkzGetPoint{prePc}
    \tkzInterLL(P,prePc)(A,B) \tkzGetPoint{Pc}
    \tkzDefLine[bisector](B,P,C) \tkzGetPoint{prePa}
    \tkzInterLL(P,prePa)(B,C) \tkzGetPoint{Pa}
    \tkzDefLine[bisector](C,P,A) \tkzGetPoint{prePb}
    \tkzInterLL(P,prePb)(C,A) \tkzGetPoint{Pb}

    \tkzLabelPoint[below left](A){$A$}
    \tkzLabelPoint[below right](B){$B$}
    \tkzLabelPoint[above](C){$C$}
    %\tkzLabelPoint[below left](P){$P$}
    \tkzLabelPoint[above right](Pa){$A_p$}
    \tkzLabelPoint[above left](Pb){$B_p$}
    \tkzLabelPoint[below](Pc){$C_p$}

    \tkzMarkAngle[arc=lll,size=1.2,mark=|||](A,P,Pc)
    \tkzMarkAngle[arc=lll,size=1.2,mark=|||](Pc,P,B)
    \tkzMarkAngle[arc=ll,size=1.2,mark=||](B,P,Pa)
    \tkzMarkAngle[arc=ll,size=1.2,mark=||](Pa,P,C)
    \tkzMarkAngle[arc=l,size=1.2,mark=|](C,P,Pb)
    \tkzMarkAngle[arc=l,size=1.2,mark=|](Pb,P,A)

    \tkzDrawSegments[line width=0.2mm](P,A P,B P,C)
    \tkzDrawSegments[line width=0.2mm,dashed](P,Pa P,Pb P,Pc)
    \tkzDrawPolygon[line width=0.3mm](A,B,C)
    \tkzDrawPoints[size=3,color=black,fill=black!50](A,B,C,P,Pc,Pb,Pa)
\end{tikzpicture}
\end{comment}
    \end{center}
    \subcaption{%
nierówność Barrowa% % lang-pl
disuguaglianza di Barrow% % lang-it
}
    \label{erdos_mordell_barrowb}
\end{minipage}
\caption{}
\end{figure}

\label{subsection_erdos_mordell}
Nierówność pojawi się w wakacyjnym wydaniu American Mathematical Monthly z 1935 roku jako problem 3740 od Erdősa. % lang-pl
La disuguaglianza apparve nel numero estivo dell'\emph{American Mathematical Monthly} del 1935 come problema 3740 proposto da Erdős. % lang-it
\index[persons]{Erdős, Paul}%
Na dowód przyjdzie poczekać. % lang-pl
Per una dimostrazione si dovrà attendere ancora qualche tempo. % lang-it
Pierwszy podadzą Louis Mordell i~David Barrow \cite{mordell_barrow_1937} w 1937 roku. % lang-pl
La prima sarà fornita da Louis Mordell e David Barrow \cite{mordell_barrow_1937} nel 1937. % lang-it
\index[persons]{Mordell, Louis Joel}%
\index[persons]{Barrow, David}%
Po dwudziestoletniej przerwie pojawi się prostszy dowód Kazarinowa \cite{kazarinoff_1957}; % lang-pl
Dopo una pausa di vent'anni comparirà una dimostrazione più semplice di Kazarinoff \cite{kazarinoff_1957}; % lang-it
\index[persons]{Kazarinoff, Donat}% Donat Konstantinovich Kazarinoff
Leon Bankoff \cite{bankoff_1958} uzasadni nierówność wykorzystując jedynie podobieństwo trójkątów oraz fakt, że $t + 1/t \ge 2$ dla $t \ge 0$; cały dowód zmieści się na połowie strony w $\Delta_{90}^8$. % lang-pl
Leon Bankoff \cite{bankoff_1958} dimostrerà la disuguaglianza utilizzando soltanto la similitudine dei triangoli e il fatto che $t + 1/t \ge 2$ per $t \ge 0$; l'intera dimostrazione occuperà mezza pagina in $\Delta_{90}^8$. % lang-it
\index[persons]{Bankoff, Leon}%
Ostatni współczesny dowód pochodzi od Claudi Alsiny i Rogera Nelsena \cite{alsina_2007}. % lang-pl
L'ultima dimostrazione moderna è dovuta a Claudi Alsina e Roger Nelsen \cite{alsina_2007}. % lang-it
\index[persons]{Alsina, Claudi}%
\index[persons]{Nelsen, Roger}%

% https://en.wikipedia.org/wiki/Erdős-Mordell_inequality
Coxeter \cite[s. 25]{coxeter_1967} przetworzy je po latach w ćwiczenie dla Czytelnika; podobnie postąpi Audin \cite[s. 102]{audin_2003}, nie szczędząc przy tym wskazówek. % lang-pl
Anni dopo Coxeter \cite[p. 25]{coxeter_1967} la trasformerà in un esercizio per il lettore; lo stesso farà Audin \cite[p. 102]{audin_2003}, senza risparmiare suggerimenti. % lang-it
Guzicki \cite[s. 380-385]{guzicki_2021} pokaże cztery dowody. % lang-pl
Guzicki \cite[p. 380-385]{guzicki_2021} presenterà quattro dimostrazioni. % lang-it
W geometrii absolutnej nierówność Erdősa-Mordella jest równoważna temu, że suma kątów trójkąta nie przekracza $\pi$ (patrz Pambuccian \cite{pambuccian_2008}). % lang-pl
Nella geometria assoluta, la disuguaglianza di Erdős-Mordell è equivalente all'affermazione che la somma degli angoli di un triangolo non supera $\pi$ (si veda Pambuccian \cite{pambuccian_2008}). % lang-it
\index[persons]{Pambuccian, Victor}%

Wzmocnieniem nierówności Erdősa-Mordella będzie nierówność Barrowa: % lang-pl
Un rafforzamento della disuguaglianza di Erdős-Mordell è la disuguaglianza di Barrow: % lang-it
\index{nierówność!Barrowa} % lang-pl
\index{disuguaglianza!di Barrow} % lang-it

% TODO: https://en.wikipedia.org/wiki/Barrow%27s_inequality

\begin{theorem}
[nierówność Barrowa] % lang-pl
[disuguaglianza di Barrow] % lang-it
    Niech $P$ będzie punktem wewnątrz trójkąta $\triangle ABC$, zaś $A_p$, $B_p$, $C_p$ punktami przecięć dwusiecznych trzech kątów wyznaczanych przez $P$ i pary wierzchołków trójkąta; tak jak na rysunku \ref{erdos_mordell_barrowb}. % lang-pl
    Sia $P$ un punto interno al triangolo $\triangle ABC$, e siano $A_p$, $B_p$, $C_p$ i punti d'intersezione delle bisettrici dei tre angoli determinati da $P$ e dalle coppie di vertici del triangolo, come nella figura \ref{erdos_mordell_barrowb}. % lang-it
    Wtedy % lang-pl
    Allora % lang-it
    \begin{equation}
        |PA| + |PB| + |PC| \ge 2 (|PA_p| + |PB_p| + |PC_p|).
    \end{equation}
\end{theorem}

Dowód Barrowa i Mordella \cite{mordell_barrow_1937} zostanie opublikowany w 1937 roku, ale nazwa ,,nierówność Barrowa'' przyjmie się nieco później, na pewno w 1961 roku Alexander Oppenheim będzie już jej używać. % lang-pl
La dimostrazione di Barrow e Mordell \cite{mordell_barrow_1937} sarà pubblicata nel 1937, ma il nome \emph{disuguaglianza di Barrow} si affermerà solo più tardi; certamente nel 1961 Alexander Oppenheim lo utilizzava già. % lang-it
\index[persons]{Oppenheim, Alexander}%
% Oppenheim, A. (1961), "New inequalities for a triangle and an internal point", Annali di Matematica Pura ed Applicata, 53: 157–163, doi:10.1007/BF02417793, MR 0124774
Mordell \cite{mordell_1962} wróci do nierówności po latach i poda prostszy dowód. % lang-pl
Mordell \cite{mordell_1962} tornerà sulla disuguaglianza dopo molti anni e ne fornirà una dimostrazione più semplice. % lang-it
Jan Ciach zaproponuje w~$\Delta_{95}^2$ (zadanie 296 w Klubie 44) inną nierówność dla punktu $P$ wewnątrz trójkąta $\triangle ABC$, w której występują odległości $R_a, R_b, R_c$ punktu $P$ od wierzchołków oraz $r_a, r_b, r_c$ od prostych $BC$, $CA$, $AB$: $8(\sqrt{R_bR_c}+r_a)(\sqrt{R_cR_a}+r_b)(\sqrt{R_aR_b}+r_c) \le 27R_aR_bR_c$. % lang-pl
Rozwiązanie ukaże się w~$\Delta_{95}^6$; równość zachodzi, gdy trójkąt jest równoboczny, a $P$ jest jego środkiem ciężkości. % lang-pl
Jan Ciach proporrà in $\Delta_{95}^2$ (problema 296 nel Klub 44) un'altra disuguaglianza per un punto $P$ interno al triangolo $\triangle ABC$, in cui compaiono le distanze $R_a, R_b, R_c$ di $P$ dai vertici e $r_a, r_b, r_c$ dalle rette $BC$, $CA$, $AB$: $8(\sqrt{R_bR_c}+r_a)(\sqrt{R_cR_a}+r_b)(\sqrt{R_aR_b}+r_c) \le 27R_aR_bR_c$. % lang-it
La soluzione apparirà in $\Delta_{95}^6$; l'uguaglianza vale quando il triangolo è equilatero e $P$ ne è il baricentro. % lang-it
\index[persons]{Mordell, Louis Joel}%

\index{nierówność!Erdősa-Mordella|)} % lang-pl
\index{disuguaglianza!di Erdős-Mordell|)} % lang-it

% % barrow tu jest

%

\subsection{Największy kwadrat w trójkącie} % lang-pl
\subsection{Il più grande quadrato in un triangolo} % lang-it

% https://en.wikipedia.org/wiki/Inscribed_square_in_a_triangle

Wewnątrz każdego trójkąta znajduje się co najmniej jeden kwadrat o maksymalnym polu. % lang-pl
Łatwo widać, że jego bok musi leżeć na którymś z boków trójkąta. % lang-pl
(Skonstruowanie go to zadanie \ref{konstrukcja_kwadratu_wpisanego}). % lang-pl
W~trójkącie równobocznym są nawet trzy takie kwadraty, każdy oparty o inną podstawę. % lang-pl

All'interno di ogni triangolo esiste almeno un quadrato di area massima. % lang-it
È facile vedere che uno dei suoi lati deve giacere su un lato del triangolo. % lang-it
(La sua costruzione è l'esercizio \ref{konstrukcja_kwadratu_wpisanego}). % lang-it
In un triangolo equilatero esistono addirittura tre quadrati di questo tipo, ciascuno costruito su una base diversa. % lang-it

\begin{proposition}
\label{trojkat_calabiego}%
    W trójkącie o dwóch ramionach długości $1$ oraz podstawie, której długość jest największym dodatnim pierwiastkiem równania % lang-pl
    In un triangolo con due lati uguali di lunghezza $1$ e base la cui lunghezza è la maggiore radice positiva dell'equazione % lang-it
    \begin{equation}
        2x^3 - 2x^2 - 3x + 2 = 0,
    \end{equation}
    można wpisać trzy kwadraty o największym możliwym polu. % lang-pl
    si possono inscrivere tre quadrati di area massima possibile. % lang-it
    $x \approx 1.55138$.
\end{proposition}

Calabi pokaże około 1997 roku, że jest to jedyny nierównoboczny trójkąt o tej własności. % lang-pl
Calabi mostrerà intorno al 1997 che questo è l'unico triangolo non equilatero con tale proprietà. % lang-it
Zadanie wpisania w trójkąt równoległoboku o największym polu przypisze się Euklidesowi; Jarosław Górnicki opisze w $\Delta_{01}^{9}$ jego rozwiązanie geometryczne (w trójkąt $ABC$ wpisujemy równoległobok $ADEF$ o jednym wierzchołku w $A$; największe pole ma ten, dla którego punkty $D$, $E$, $F$ są środkami boków) i porówna je z rozwiązaniem analitycznym, które korzysta z pochodnej. % lang-pl
Il problema di inscrivere in un triangolo un parallelogramma di area massima sarà attribuito a Euclide; Jarosław Górnicki descriverà in $\Delta_{01}^{9}$ la sua soluzione geometrica (nel triangolo $ABC$ si inscrive un parallelogramma $ADEF$ con un vertice in $A$; l'area massima si ottiene quando i punti $D$, $E$, $F$ sono i punti medi dei lati) e la confronterà con la soluzione analitica, che usa la derivata. % lang-it
% Calabi, E. "Outline of Proof Regarding Squares Wedged in Triangle." Nov. 3, 1997. http://algo.inria.fr/csolve/calabi.html.

%

% https://www.awesomemath.org/wp-pdf-files/math-reflections/mr-2025-04/mr_4_2025_trigonometric_inequalities_v2.pdf

% TODO: https://en.wikipedia.org/wiki/List_of_triangle_inequalities#

%